\section{A gallery of counterexamples}
\label{sec:gallery}

Twelve functions carry the whole argument. Each entry gives the definition,
the property the function has, the property it lacks, and the implication it
refutes.

\begin{center}
\begin{tabular}{@{}lll@{}}
\toprule
& \textbf{function} & \textbf{refutes} \\
\midrule
$\mathrm{E}1$  & removable spike      & limit $\Rightarrow$ continuous \\
$\mathrm{E}2$  & $|x|$                & continuous $\Rightarrow$ differentiable \\
$\mathrm{E}3$  & $x^2\sin(1/x)$       & differentiable $\Rightarrow$ $\mathscr{C}^1$ \\
$\mathrm{E}4$  & $\sin(1/x)$          & bounded $\Rightarrow$ has a limit \\
$\mathrm{E}5$  & step function        & integrable $\Rightarrow$ continuous \\
$\mathrm{E}6$  & Thomae               & integrable $\Rightarrow$ disc.\ set nowhere dense \\
$\mathrm{E}7$  & Dirichlet            & bounded $\Rightarrow$ integrable \\
$\mathrm{E}8$  & Weierstrass          & continuous $\Rightarrow$ differentiable somewhere \\
$\mathrm{E}9$  & $x^n$ on $[0,1]$     & pointwise preserves continuity \\
$\mathrm{E}10$ & $\sin(nx)/\sqrt n$   & uniform conv.\ commutes with $d/dx$ \\
$\mathrm{E}11$ & tall spikes          & pointwise preserves $\int$ \\
$\mathrm{E}12$ & Volterra             & derivative $\Rightarrow$ integrable \\
\bottomrule
\end{tabular}
\end{center}

\medskip
\noindent
$\mathrm{E}1$ \emph{Removable spike.} $f(x)=0$ for $x\ne0$ and $f(0)=1$. The
limit at $0$ is $0$ and the value is $1$.

\smallskip\noindent
$\mathrm{E}2$ \emph{Absolute value.} $f(x)=|x|$ is continuous everywhere; at
$0$ the difference quotient $|h|/h$ equals $\pm1$ according to the side, so
$f'(0)$ does not exist. Note $f$ is Lipschitz, so that condition does not
give differentiability either.

\smallskip\noindent
$\mathrm{E}3$ \emph{$x^2\sin(1/x)$.} With $f(0)=0$, the bound $|f(x)|\le x^2$
gives $f'(0)=0$, while for $x\ne0$
\[ f'(x) = 2x\sin\tfrac1x - \cos\tfrac1x, \]
which has no limit at $0$. Thus $f$ is differentiable everywhere with $f'$
discontinuous at $0$, the discontinuity being of the second kind as
Corollary~\ref{cor:nojump} requires. This function is the building block of
$\mathrm{E}12$.

\smallskip\noindent
$\mathrm{E}4$ \emph{$\sin(1/x)$.} Bounded by $1$ with no limit at $0$: along
$x_n=1/(2n\pi)$ the values are $0$ and along $y_n = 1/(2n\pi+\pi/2)$ they are
$1$. Its oscillation at $0$ is $2$.

\smallskip\noindent
$\mathrm{E}5$ \emph{Step function.} $f=0$ on $[0,1)$ and $f=1$ on $[1,2]$.
Integrable, having one discontinuity, and not continuous; by
Corollary~\ref{cor:nojump} it is also not a derivative.

\smallskip\noindent
$\mathrm{E}6$ \emph{Thomae's function.} $f(p/q)=1/q$ in lowest terms and
$f=0$ at irrationals. Continuous at every irrational, discontinuous at every
rational, and integrable with $\int_0^1 f = 0$ by
Theorem~\ref{thm:lebesgue}, $\Q$ being countable. No function is continuous
exactly at the rationals, so $\mathrm{E}6$ has no mirror image
\cite[appendix to Chapter~4]{annotated}.

\smallskip\noindent
$\mathrm{E}7$ \emph{Dirichlet's function.} $\1_\Q$ is bounded with
$\omega_f\equiv1$, so it is nowhere continuous and not integrable; directly,
every upper sum is $1$ and every lower sum $0$.

\smallskip\noindent
$\mathrm{E}8$ \emph{Weierstrass's function.}
$f(x)=\sum_{k\ge0} a^k\cos(b^k\pi x)$ with $0<a<1$, $b$ an odd integer, and
$ab>1+\tfrac32\pi$. The series converges uniformly, so $f$ is continuous by
Theorem~\ref{thm:unifcont}, and $f$ is differentiable at no point. Its
partial sums are $\mathscr{C}^\infty$ and converge uniformly, so
$\mathrm{E}8$ also witnesses that uniform convergence does not preserve
differentiability (Theorem~\ref{thm:unifdiff}). By Baire's
theorem the continuous functions differentiable somewhere are of first
category in $\mathscr{C}[a,b]$, so this behaviour is typical rather than
exceptional.

\smallskip\noindent
$\mathrm{E}9$ \emph{$x^n$ on $[0,1]$.} Each term is continuous and the
pointwise limit is not; the convergence is not uniform, since
$\sup_{[0,1]}|f_n-f| = 1$ for every $n$.

\smallskip\noindent
$\mathrm{E}10$ \emph{$\sin(nx)/\sqrt n$.} Converges uniformly to $0$ while
the derivatives $\sqrt n\cos(nx)$ diverge everywhere, so the interchange
$\lim f_n' = (\lim f_n)'$ fails under uniform convergence. Differentiation
multiplies amplitude by frequency, and uniform convergence bounds only
amplitude.

\smallskip\noindent
$\mathrm{E}11$ \emph{Tall spikes.} $f_n = n\1_{(0,1/n)}$, each integrable
with $\int_0^1 f_n = 1$, tending pointwise to $0$. The sequence admits no
integrable dominating function --- near $0$ the envelope $\sup_n f_n$ grows
like $1/x$ --- which is exactly what dominated convergence would require.

\smallskip\noindent
$\mathrm{E}12$ \emph{Volterra's derivative.} The derivative $F'$ of the
function constructed in \S\ref{sec:int}: bounded, a derivative by
construction, and not Riemann integrable, because it oscillates on a nowhere
dense set of positive measure. The function $F$ itself is continuous, hence
integrable; it is the derivative that witnesses the gap.

\section{Concluding remarks}

Three observations summarise the map.

The local chain is strict at every link and the witnesses are elementary:
$\mathrm{E}1$, $\mathrm{E}2$, $\mathrm{E}3$ suffice to separate having a
limit, continuity, differentiability, and continuous differentiability.

Integrability is not part of that chain. It is global, it follows from
compactness rather than from any pointwise hypothesis
(Theorem~\ref{thm:ci}), and it admits an exact characterisation in terms of
the size of the discontinuity set (Theorem~\ref{thm:lebesgue}), for which
oscillation is the natural language.

The interesting failures are all failures of interchange: of $\lim$ with
$\int$, of $\lim$ with $\lim$, and of $\lim$ with $d/dx$. Uniform convergence
repairs the first two and cannot repair the third, because differentiation
amplifies precisely what uniformity controls. Each gap in the lattice invites
the question of what hypothesis would close it: pointwise convergence
strengthened to uniform restores continuity and the integral; uniform
convergence transferred to the derivatives restores differentiation
(Theorem~\ref{thm:repair}); and replacing the Riemann integral by the
Lebesgue integral restores the interchange under a domination hypothesis in
place of a uniformity one.
